\section{Literature}
\label{sec:literature}

Early related work is the classical economics of collusion. \citet{green1984} study repeated competition where firms can only observe each other's market shares with noise, not their actions directly, and show that prices above the competitive level can be supported in equilibrium, but only by occasional price wars that get triggered when demand looks suspiciously low. The reason is that firms cannot tell whether a rival cheated on the implicit agreement or whether demand simply fell, so any sufficiently bad signal has to trigger a punishment to keep cheating off the table. This is the textbook explanation of why markets like petrol retailing show recurring price war episodes even when the firms involved seem to be coordinating most of the time. \citet{maskin1988} build a different repeated pricing model in which firms take turns setting prices, and find that this alternating move structure can sustain either price cycles (Edgeworth cycles) or stable high prices. The two papers introduce the basic ideas of imperfect monitoring, alternating moves, and coordination on a focal high price, which the recent algorithmic collusion literature reuses.

\citet{calvano2020} provide the foundational result for this thesis. They place two independent tabular Q-learning agents \citep{watkins1992} in a repeated Bertrand pricing game with logit demand and a discrete grid of prices. Each agent keeps a Q-table that records what each action has been worth so far in each state, updates the table after every period using the standard Q-learning rule with $\varepsilon$-greedy exploration, and earns the period's profit as its reward. The two agents never communicate. Yet they reliably converge on prices well above the competitive Nash level, typically reaching about 84\% of the way from Nash to joint monopoly. The result holds across a wide range of hyperparameter choices and exploration rules. The strategies the agents end up with also seem to behave like reward punishment schemes. When one agent deviates by undercutting, the other temporarily lowers its price too (a short price war), and once the deviation ends the prices return to the high level. So in welfare terms the outcome looks like a cartel, but without any of the illegal communication that prosecuting a cartel would normally require.

A follow-up literature has started to relax the strong assumptions of the Calvano environment one at a time. \citet{calvano2021} look at imperfect monitoring through a stochastic demand shock in a Cournot setting and find that algorithmic collusion shrinks but does not disappear at moderate noise levels. \citet{klein2021} studies sequential price-setting as the base game (rather than as a deviation from the simultaneous one), and finds that Q-learners under alternating moves also converge on supracompetitive prices, sometimes through the kind of cyclic dynamics Edgeworth described. \citet{calvano2023} then ask whether the high prices Q-learners settle on are sustained by real punishment threats or are just side effects of the rule the algorithm uses to pick actions, and propose off-path deviation tests as a way to tell the two apart. The test manually forces one agent to deviate after convergence and watches whether the other retaliates. \citet{asker2024} vary the learning algorithm itself rather than the information environment, comparing synchronous and asynchronous Q-learning protocols, and show that the choice of algorithm matters a lot for the prices that emerge.

Those papers also give benchmarks for interpreting the single-friction results in this thesis. The noisy-monitoring result below is consistent with \citet{calvano2021} and \citet{zhang2025}: imperfect rival information weakens learned collusion, and sufficiently strong noise can push the outcome close to the competitive benchmark. The asynchronous-pricing result is more mixed relative to \citet{klein2021}. Like Klein, I find that alternating moves do not eliminate supracompetitive pricing. Unlike Klein's sequential-pricing environment, however, imposing alternation on the Calvano baseline lowers the collusion score sharply before it settles on a lower plateau. This difference is useful because it suggests that asynchrony has two roles: it can support coordination in a game designed around sequential moves, while also slowing or disrupting learning when added as a friction to an otherwise simultaneous game. The off-path tests follow \citet{calvano2023} by asking whether the resulting high prices are backed by punishment behaviour rather than mere convergence inertia.

There is also early empirical evidence from real markets. \citet{assad2024} use station level data on the German retail gasoline market and find that markups rose around the time the major chains adopted algorithmic pricing software. Their difference in differences design uses chain headquarters' adoption decisions as the source of variation. \citet{brown2023} use high-frequency online retail data and document that firms running pricing algorithms at different update frequencies sustain higher equilibrium prices than firms updating at the same speed, a finding that links the asynchronous pricing literature directly with the algorithmic collusion one. Together, these two papers suggest that the Calvano simulation is not just a numerical exercise. It captures something that is already happening in markets where algorithmic pricing is widespread.

On the policy side, \citet{oecd2025} surveys competition policy across the G7 and notes both the underlying problem (current law relies on detecting communication, which algorithmic agents do not need) and the early enforcement responses to specific cases. \citet{zhang2025} asks whether mandated noise injection on the rival's observed price could be a regulatory tool, and finds that adding enough bounded Laplacian noise can push learned prices back to the competitive level in a single friction setting.

What is missing from all of this is any controlled study of what happens when several frictions are imposed at the same time. The one-friction literature is informative, and the single-friction experiments in this thesis deliberately connect back to it: noise weakens collusion, alternating moves still leave positive collusion, and off-path tests can diagnose whether the learned outcome is strategic. But real markets typically have several of these frictions going on at once. A regulator considering more than one policy instrument needs to know whether the instruments add up, cancel out, or interact in some other way. Combining a delayed disclosure rule with a noise injection rule, for example, could disrupt collusion more than either one alone, less than either alone, or something in between. The simulation literature does not currently answer that. This thesis is an attempt to start, using a single controlled environment in which the three frictions are introduced one at a time and then jointly in a combined-friction design.
